A nearest-neighbour interaction on $\Z^2$ is \emph{computably presented}
if its finite alphabet is supplied and one algorithm outputs, for each lattice
edge $\{x,y\}$, the complete table
$b_{\{x,y\}}:S^2\to\mathbb Q_{>0}$. The tables may depend on absolute position.
We use this as a promise class: an index is promised to describe a total
algorithm with positive entries, just as in Proposition~\ref{prop:phase-arith}.
No decision procedure for validity of arbitrary program indices is asserted.
All reductions below output uniformly total valid presentations on every input.
Completeness is under computable many-one reductions within this promise class.

\begin{theorem}[Sharp bounds for position-dependent presentations]\label{thm:cp-sharp}
For computably presented nearest-neighbour interactions on $\Z^2$ at $\beta=1$:
\begin{enumerate}[label=(\roman*)]
\item For a fixed decidable $R$ defining a $\Sigma^0_2$-complete set, there is a
computable map $e\mapsto\Phi_e$, with $S=\{-1,+1\}$ and edge weights in
$\{1,16,256\}$, such that
\[
 \#\operatorname{Ext}\G(\Phi_e)=
 \begin{cases}2,&\exists w\,\forall n\,R(e,w,n),\\1,&\text{otherwise}.\end{cases}
\]
Thus $\mathsf{NONUNIQUE}$ is $\Sigma^0_2$-complete and $\mathsf{UNIQUE}$ is
$\Pi^0_2$-complete.
\item For every fixed $r\ge2$, exactly $r$ extremal Gibbs states is
$D_2(\Sigma^0_2)$-complete, and at least $r$ is $\Sigma^0_2$-complete,
already on product-Potts, respectively Potts, nearest-neighbour interactions
with $r(r+1)$, respectively $r$, on-site states and positive integer weights.
\item On every unique input of (i), and on every unique input of the family $\Phi_C$ in (iv), the unique Gibbs state satisfies
\[
 |\Cov_\mu(f,g)|\le C_0|A_f||A_g|\|f\|_\infty\|g\|_\infty
 e^{-a_0\,d_1(A_f,A_g)}
\]
for all bounded local $f,g$, with fixed $C_0,a_0>0$ independent of $e$.
Every local $f$ has finite-volume boundary oscillation
$\Delta_f(L)=\max_\eta\gamma^\eta_{B_L}(f)-\min_\eta\gamma^\eta_{B_L}(f)=0$,
over all collar configurations $\eta$, for all sufficiently large $L$.
\item Let $\Phi_C$ be the binary square-boundary family constructed in the
proof, with $C\subseteq\mathbb N_{\ge1}$ given by a promised total decider.
Write $m^*=\nu^+(\sigma_0)>0$ for the homogeneous $16$-weight Ising magnetization
and $f_0=\one_{\{\sigma_0=+1\}}$. For each fixed rational $0<c\le m^*$
(in particular $c=1/4$), no total computable function bounds
$\min\{L\ge0:\Delta_{f_0}(L)<c\}$ on all unique inputs.
On a nonunique input every observable inside the outermost selected square
has zero Gibbs oscillation, while single-spin indicators sufficiently far
outside distinguish the two extremal states with oscillation at least $m^*/2$.
No total computable function bounds the distance from the origin needed to
find such a spin, even on the family of (i).
\end{enumerate}
\end{theorem}
\begin{proof}
\emph{The square-boundary selector.}
Fix an integer $q\ge2$ and an integer $M>1+\sqrt q$. Put
$B_s=[-s,s]^2\cap\Z^2$, and let $D_s$ be its nearest-neighbour edge boundary.
For a decidable $C\subseteq\mathbb N_{\ge1}$ use the factors
\begin{equation}\label{eq:selector-weights}
 b_{\{x,y\}}(u,v)=
 \begin{cases}
 M^{\one_{\{u=1\}}+\one_{\{v=1\}}},&\{x,y\}\in\bigcup_{s\in C}D_s,\\
 M^{\one_{\{u=v\}}},&\text{otherwise}.
 \end{cases}
\end{equation}
Each edge crosses at most one $D_s$: its smaller endpoint sup norm is the
only possible $s$. One query to the decider for $C$ therefore computes its
table. Every weight belongs to $\{1,M,M^2\}$. The first line factors into
two one-site factors, so it cuts the interaction between its endpoints and
gives each endpoint the usual colour-$1$ boundary factor on that missing edge.
The binary version identifies colour $1$ with $+1$.

If $C$ is infinite, all components $V_i$ of the ordinary-edge graph are finite.
Let $\pi_{V_i}^1$ be the homogeneous Potts law on $V_i$ with colour-$1$
boundary conditions on every missing edge, including both sides of an annulus.
The kernel on an entire component is independent of its exterior. The DLR
equation on any finite union of components consequently forces its law to be
the corresponding product. Cylinders determine a measure, giving
\begin{equation}\label{eq:selector-product}
 \G(\Phi_C)=\left\{\bigotimes_i\pi_{V_i}^1\right\}.
\end{equation}
This product is Gibbs: given a finite $\Lambda$, include all components it
meets in a finite union $U$. Its law is invariant under $\gamma_U$, and
consistency gives $\gamma_U\gamma_\Lambda=\gamma_U$.
Thus infinitely many components, regardless of their sizes, leave no
unspecified cylinder marginal.

If $C$ is finite, only finitely many edge factors change. For the homogeneous
specification define the positive cylinder function
\[
 F_C(\sigma)=\prod_{\{x,y\}\in\bigcup_{s\in C}D_s}
 M^{\one_{\{\sigma_x=1\}}+\one_{\{\sigma_y=1\}}
       -\one_{\{\sigma_x=\sigma_y\}}}.
\]
For every finite $\Lambda$ the modified conditional expectation is
$\gamma'_\Lambda(f)=\gamma_\Lambda(fF_C)/\gamma_\Lambda(F_C)$;
factors outside $\Lambda$ cancel. Hence
\[
 T_C\nu=F_C\nu/\nu(F_C)
\]
is Gibbs for the modified specification whenever $\nu$ is Gibbs for the
homogeneous one. Indeed, setting
$h=\gamma_\Lambda(fF_C)/\gamma_\Lambda(F_C)$, which is exterior measurable,
gives $\nu(F_Ch)=\nu(\gamma_\Lambda(F_Ch))=\nu(fF_C)$.
Reweighting by $F_C^{-1}$ is the inverse. Both functions are bounded and
bounded away from zero. Mutual absolute continuity transfers tail triviality,
and hence extremality \cite[Theorem 7.7]{Georgii}; this direct proof does not
require the normalized reweighting to be affine (cf.\ \cite[Section 7.4]{Georgii}).
Since $M>1+\sqrt q$, that is $\log M>\beta_c(q)=\log(1+\sqrt q)$, the classification of all DLR states of the homogeneous planar Potts model gives exactly $q$ extremal states, its pure phases \cite[Theorem 1.1 and Corollary 1.3]{CDIV2014}.
For $q=2$ this is the Aizenman--Higuchi classification~\cite{Aizenman1980,Higuchi1981}.
Thus the selector has exactly $q$ extremal states for finite $C$, and one for
infinite $C$. The set $C$ is allowed to be empty, giving the homogeneous model.
Simply deleting the cut edges would already give this hardness construction;
the positive colour-$1$ factors are included for the clustering assertion below.

\emph{Compiling the quantifiers.}
For (i) put
\[
 h_e(s)=\min\bigl(\{w\le s:\forall n\le s\ R(e,w,n)\}\cup\{s+1\}\bigr),
 \qquad C_e=\{s\ge1:h_e(s)>h_e(s-1)\}.
\]
These are uniformly computable. A refuted candidate never returns, and a newly
admitted candidate cannot be smaller than the preceding minimum, so $h_e$ is
nondecreasing. If a successful witness exists, the least one eventually
survives after all smaller candidates are refuted, so $h_e$ stabilizes.
Conversely, stabilization at $w$ implies $R(e,w,n)$ for every $n$.
If there is no witness, every fixed candidate is eventually refuted and
$h_e(s)\to\infty$. Thus $C_e$ is finite exactly on the stated $\Sigma^0_2$ set.
Take $q=2$, $M=16$ in the selector. Proposition~\ref{prop:phase-arith} supplies
the matching upper bounds. No validity promise is needed by this reduction.

\emph{Products and exact counts.}
For (ii), apply the selector to a $q=r$ layer for an arbitrary
$A\in\Sigma^0_2$ and a $q=r+1$ layer for an arbitrary $B\in\Sigma^0_2$.
Choose fixed integers $M_q>1+\sqrt q$ for these two layers. Multiplying their
edge factors gives positive integer tables on the product alphabet.
We justify both directions of multiplication of extremal counts.

First, any joint Gibbs law is invariant under the finite-volume kernel
resampling only one layer: composing that kernel after the product DLR
kernel leaves the latter unchanged. Testing against arbitrary functions of
the other layer shows that the conditional law of the first layer given the
entire second layer is almost surely Gibbs. If the first marginal is extreme,
these conditional laws have that extreme barycenter and are almost surely
equal to it, by uniqueness of the extreme decomposition
\cite[Theorem 7.26]{Georgii}. The joint law is then the product of its marginals.
An extreme joint law has tail-trivial marginals, since each layer's tail is
contained in the joint tail; its marginals are Gibbs and extreme
\cite[Theorem 7.7]{Georgii}. It is therefore a product of extreme marginals.

Conversely, a product of two extreme Gibbs laws is Gibbs. In any convex
decomposition of that product into joint Gibbs laws, extremality of each
marginal forces both summands to have the same fixed marginals. The preceding
conditional-law argument then forces each summand to be the same product.
The product is extreme. This proves the converse without identifying the
joint tail with a product of tail sigma-fields.
The possible counts are consequently $1,r,r+1,r(r+1)$ according as neither,
only $A$, only $B$, or both predicates hold. The count is exactly $r$ precisely
on $A\setminus B$, proving $D_2(\Sigma^0_2)$ hardness. A single $q=r$ layer
proves the at-least-$r$ hardness. Proposition~\ref{prop:phase-arith} supplies
the upper bounds. At $r=1$ exact count is uniqueness, which is
$\Pi^0_2$-complete by (i); no sharpness claim for infinitely many states is made.

\emph{Uniform clustering.}
For the binary model $M=16$ corresponds to Ising coupling $K=\frac12\log16$.
For any finite $V\subset\Z^2$, including multiply connected sets, the
plus-boundary law obeys
\begin{equation}\label{eq:selector-covariance}
 0\le\Cov_{\pi_V^+}(\sigma_x,\sigma_y)
 \le\Cov_{\nu^+}(\sigma_x,\sigma_y)
 \le C_0e^{-a_0|x-y|_1},
\end{equation}
with constants depending only on $K$; see
\cite[Section 2.3.1, (2.32)--(2.33)]{BCK2003} and the primary decay result
\cite{CCS1987}. The middle inequality follows from GHS~\cite{GHS1970}:
in a larger plus-boundary box pin the spins outside $V$ by increasing their
nonnegative fields to infinity, which decreases truncated two-spin
correlations; then send the larger box to infinity. On the unique branch,
\eqref{eq:selector-product} gives zero covariance between different components
and the same upper bound within a component.

For completeness, the extension to arbitrary real local observables has a
short FKG proof (compare the association inequalities in~\cite{BulinskiShashkin2007}).
Let $\delta_x(f)$ be the maximum change in $f$ on flipping only spin $x$, and
put $L_f=\frac12\sum_x\delta_x(f)\sigma_x$.
Both $L_f+f$ and $L_f-f$ are increasing. The component laws, and their product,
are associated. Summing the FKG inequalities for the two opposite-sign pairs
gives $\Cov(f,g)\le\Cov(L_f,L_g)$; the two same-sign pairs give the reverse
bound. Therefore
\[
 |\Cov(f,g)|\le\frac14\sum_{x,y}\delta_x(f)\delta_y(g)
                   \Cov(\sigma_x,\sigma_y).
\]
Since $\delta_x(f)\le2\|f\|_\infty$, \eqref{eq:selector-covariance} proves
(iii).
If $C$ is infinite and $f$ is supported in $B_s$ for a selected $s$, all
components meeting its support lie inside $B_s$. Thus the explicit forgetting
radius
\[
 s_*(f,C)=\min\{s\in C:\supp f\subseteq B_s\}
\]
satisfies $\Delta_f(L)=0$ for every $L\ge s_*(f,C)$. Here volumes are $B_L$
and maxima and minima range over their full collars.

\emph{No total scale or location bound.}
For (iv), work on the uniformly decidable family of cut sets, rather than
assuming an arbitrary fixed relation $R$ contains every halting instance.
Given a machine $p$, let $C_p=\{s\ge1:p\text{ halts within }s\text{ steps}\}$,
with the machine convention arranged so a halt takes at least one step.
This is a uniformly decidable set. If $p$ never halts it is empty and the
model is homogeneous and nonunique. If $p$ halts at time $t$, it is
$\{t,t+1,\ldots\}$ and the model is unique. For every $L<t$, the origin kernel
is still the homogeneous kernel. Its plus/minus boundary indicator
oscillation is at least $m^*$, by monotonicity of finite-volume plus
magnetization. For an explicit rational threshold, the elementary Peierls
contour bound gives
\[
 \nu^+(\sigma_0=-1)\le\sum_{\ell\ge4}4\ell(3/16)^\ell<1/4,
 \qquad m^*>1/2,
\]
so $c=1/4$ is valid. The same argument works for every fixed rational
$0<c\le m^*$. A total upper bound $b(p)$ on the first successful $L$ on unique
inputs would imply $t\le b(p)$ whenever $p$ halts, and simulation for $b(p)$
steps would decide halting. This proves the scale obstruction. Quantum
system-size thresholds beyond computable bounds are a related prior
phenomenon~\cite{BCLPW2018}; the present forgetting is exact after a finite cut.

If $C$ is finite with last selected radius $s$, every Gibbs law has the same
marginal on $B_s$. Far from the finite support of $F_C$, exponential clustering
of the homogeneous plus and minus phases gives
\[
 T_C\nu^+(\sigma_x)\longrightarrow m^*,\qquad
 T_C\nu^-(\sigma_x)\longrightarrow -m^*
 \quad (|x|_1\to\infty).
\]
Indeed the difference from the unmodified expectation is
$\Cov_{\nu^\pm}(\sigma_x,F_C)/\nu^\pm(F_C)$, which tends to zero by the
same FKG estimate and spin-flip symmetry. Single-spin indicator oscillations
therefore tend to $m^*$. When $C$ is empty this holds at every site; the
statement about an outermost selected square is then vacuous.
If a total computable radius bounded a site of oscillation at least $m^*/2$
on every nonunique input of (i), finitely many indicator searches in that
box, each below $1/4<m^*/2$, would jointly semidecide uniqueness. On unique
inputs all converge to zero; on nonunique inputs at least one can never pass.
This contradicts the $\Pi^0_2$ completeness in (i). It also rules out such a
bound on the larger uniformly decidable-$C$ family.
\end{proof}

Giving every edge outside one coordinate plane the constant table $1$ embeds these lower bounds in $\Z^d$ for every $d\ge2$: the Gibbs states are the products of those of the plane model with independent uniform spins elsewhere.

The computation in this construction sits in the dependence of the tables on
absolute position. Each increase of the least surviving candidate places a
plus-boundary cut at that stage's radius. This is a position-dependent
construction and supplies no translation-invariant interaction realizing
the same hierarchy completeness.

\begin{conjecture}\label{conj:ti-sharp}
$\mathsf{NONUNIQUE}$ is $\Sigma^0_2$-complete on translation-invariant
finite-range interactions on $\Z^2$ with strictly positive rational weights
at $\beta=1$.
\end{conjecture}

\begin{remark}[Restrictions on complete families and on translation invariance]\label{rem:hierarchy}
For a translation-invariant specification, $\G(\Phi)$ is closed under
translations, so translating a local observable preserves its Gibbs
oscillation. Suppose total computable $D(e)$ and $\delta(e)>0$ guarantee,
on every nonunique input, an observable of diameter at most $D(e)$, sup norm
at most $1$, and oscillation at least $\delta(e)$. Translate its support into
$[0,D(e)]^2$. Expanding it in the cylinder indicators of this box gives some
indicator with gap at least $\delta(e)/|S_e|^{(D(e)+1)^2}$.
The maximum of their finitely many finite-volume oscillations decreases to
zero on unique inputs and stays above that computable gap on nonunique ones.
Remark~\ref{rem:gap-ce} makes uniqueness c.e. Such a pair $(D,\delta)$ cannot
exist in a $\Sigma^0_2$-complete coexistence family. Only the computable pair is excluded: this argument does not exclude a computable diameter bound with an unspecified positive gap, or a computable gap at unbounded diameter.
Theorem~\ref{thm:cp-sharp} escapes by location, which translation invariance
forbids. T\"orm\"a's translation-invariant deterministic sparsification gives
$\Pi^0_2$-complete unique ergodicity~\cite[Proposition 48, arXiv numbering]{Torma2015};
an analogue controlling all positive-temperature DLR states remains missing.
The following further facts distinguish general restrictions from the two
that specifically use translation invariance.

\emph{(a) Every class: boundary-rate catalogues and certificates.}
Let $e\mapsto\Phi_e$ be any uniformly computable family in the effective
finite-alphabet, positive, known-range setting, translation invariant or not,
whose nonuniqueness set is $\Sigma^0_2$-complete. With a computable range $r_e\ge1$,
define
\[
 V_e(n,t)=\max_{\eta,\zeta}
 \bigl\|(\gamma_{B_{n+r_e+t}}^\eta)_{B_n}
        -(\gamma_{B_{n+r_e+t}}^\zeta)_{B_n}\bigr\|_{\TV}.
\]
These quantities are computable rationals for rational tables and computable
reals for computable-real names. For fixed $n$, their limits are the maximum
TV diameters of Gibbs marginals, by the finite-volume compactness argument
of Remark~\ref{lem:extrema}. For any total computable catalogue of nonnegative
bounds $b(e,k,n,t)$ tending to zero in $t$ for each $e,k,n$, the property
\[
 \exists k\ \forall n,t:\quad V_e(n,t)\le b(e,k,n,t)
\]
is $\Sigma^0_2$ and implies uniqueness. For computable-real quantities, the
non-strict inequality is a universal predicate on rational approximations;
its precision variable joins the universal quantifiers. If the catalogue
covered every unique output, the $\Pi^0_2$-complete uniqueness set would also
be $\Sigma^0_2$, contradicting the strict arithmetical hierarchy. Likewise,
a sound c.e. uniqueness-certificate scheme covering every unique output
would make uniqueness c.e., again impossible. These arguments hold for
\emph{every class}; neither uses translation invariance.

In particular, a bound $V_e(n,t)\le C_e|B_n|e^{-a_et}$ with arbitrary positive
constants is covered by the computable catalogue
$K|B_n|2^{-\lfloor t/K\rfloor}$, $K\ge1$: choose
$K\ge\max\{1,C_e,\log(2)/a_e\}$, and conversely each catalogue member is an
exponential envelope. Polynomial prefactors are covered by a further integer parameter: the weak-mixing condition of~\cite{MOS1994} with $\Lambda=B_{n+r_e+t}$ and $\Delta=B_n$ gives $V_e(n,t)\le C|B_n|\,|\partial^+_{r_e}B_{n+r_e+t}|\,e^{-\gamma t}$, which lies below $K|B_n|(n+t+K)^2\,2^{-\lfloor t/K\rfloor}$ for every integer $K\ge\max\{4Cr_e,\,3r_e+1,\,\log(2)/\gamma\}$. Thus some unique outputs fail the
boundary weak-mixing condition of Martinelli--Olivieri--Schonmann~\cite{MOS1994},
in its uniform cube form, in every complete family. Uniform infinite-volume
bulk clustering can nevertheless persist, as Theorem~\ref{thm:cp-sharp}(iii)
shows. The catalogue statement has quantifiers ``for each catalogue, some
unique output is missed''; it does not posit one model missed by every
computable rate. For each individual computable unique model, searching
finite-volume diameters below any prescribed accuracy supplies a convergence
modulus, without a uniform total catalogue across all inputs.

\emph{(b) Translation invariance: recoverable alignment.}
Write $\Omega=S^{\Z^d}$, $(T_z\omega)_x=\omega_{x+z}$ and $T_z\mu=\mu\circ T_z^{-1}$,
and let $\Gamma\subsetneq\Z^d$ have finite index. No map
$q:\Omega\to\Z^d/\Gamma$ can satisfy $q(T_z\omega)=q(\omega)+z$ at every
configuration, since a constant configuration is fixed by every translation.
Let therefore $\Omega_0\subseteq\Omega$ be measurable and invariant under all
translations and all finite modifications, and let $q:\Omega_0\to\Z^d/\Gamma$
be measurable, unchanged by every finite modification inside $\Omega_0$, and
equivariant, $q(T_z\omega)=q(\omega)+z$ on $\Omega_0$. Then $\Omega_0$ and each
event $\{q=a\}$ are tail measurable: for each finite region, membership and the
value of $q$ can be read after replacing that region by a fixed configuration.
Suppose a translation-covariant specification has a Gibbs state charging
$\Omega_0$. By the extreme decomposition and tail triviality of extremal states \cite[Theorems 7.26 and 7.7]{Georgii}, some
extremal state $\mu$ has $\mu(\Omega_0)=1$, and tail triviality makes $q$
almost surely equal to some $a$. The translates $T_z\mu$ are extremal and give
probability one to the respective events $\{q=a+z\}$; distinct cosets require
distinct states, so there are at least $[\Z^d:\Gamma]$ extremal states.
A unique Gibbs state is translation invariant and tail trivial, so it gives
$\Omega_0$ probability zero. Thus a translation-invariant unique branch has no
recoverable alignment on an event of positive probability, even if thermal
defects leave a decoded alignment well defined.

\emph{(c) Translation invariance: exact finite phase counts.}
If there are exactly $r$ extremal states, translations permute these $r$
states. Every permutation of $r$ objects has order dividing
$\ell_r=\operatorname{lcm}(1,\ldots,r)$. Each generator $T_{e_i}^{\ell_r}$ therefore
fixes every extremal state; their compositions fix them under all
$\ell_r\Z^d$ translations. The Gibbs simplex with these finitely many extreme
points consists of their convex combinations, so every Gibbs state is
$\ell_r\Z^d$-periodic as well. In particular, an extremal state with infinitely
many distinct translates precludes every finite exact phase count.
\end{remark}
